% Clase 2 --- Barra con densidad lineal uniforme: geometría del cálculo de la
% fuerza sobre q0 situada en el plano bisector, a distancia D del centro.
\begin{center}
\begin{tikzpicture}
  % barra
  \draw[figline,line width=1.6pt] (-2.6,0) -- (2.6,0);
  \foreach \x in {-2.4,-1.8,...,2.4} \node[figlbl,figred,scale=0.8] at (\x,0.28){$+$};
  \node[figlbl,figblue] at (2.95,0){$\lambda$};
  \node[figlblaux,below] at (-2.6,-0.05){$-L/2$};
  \node[figlblaux,below] at (2.6,-0.05){$L/2$};
  % elemento dq
  \draw[figred,line width=2pt] (1.15,0) -- (1.45,0);
  \draw[figaux] (1.3,0) -- (1.3,0.75);
  \node[figlbl,figred,above] at (1.3,0.7){$dq=\lambda\,dx$};
  % carga q0
  \node[figpos] at (0,-2.1) {};
  \node[figlbl,left=5pt] at (0,-2.1) {$q_0$};
  % distancia D
  \draw[figaux] (0,0) -- (0,-2.1) node[figlblaux,midway,left=2pt]{$D$};
  % distancia r
  \draw[figgray,line width=0.7pt] (1.3,0) -- (0,-2.1);
  \node[figlblaux] at (0.87,-1.0){$r$};
  % dF sobre q0 (repulsiva, alejándose del elemento)
  \draw[figvecres] (0,-2.1) -- (-0.6,-3.08);
  \node[figlbl,figred,below] at (-0.6,-3.05){$d\vec F$};
\end{tikzpicture}\\[0.4em]
{\small\itshape Elemento $dq=\lambda\,dx$ a distancia $r=\sqrt{x^2+D^2}$ de
$q_0$. Por simetría las componentes verticales se cancelan y sobrevive
$dF_x=dF\,\cos\theta$ con $\cos\theta=D/r$.}
\end{center}
